\section*{Streszczenie}
% \addcontentsline{toc}{chapter}{Streszczenie}
\label{chap:Streszczenie}
%\addcontentsline{toc}{chapter}{Streszczenie}
%\markboth{STRESZCZENIE}{STRESZCZENIE}

% \quad Wraz z rosnącą liczbą kierunków studiów oferowanych przez uczelnie, kandydaci coraz częściej mają trudności z podjęciem dobrze poinformowanej decyzji dotyczącej wyboru kierunku. Aby rozwiązać ten problem, autorka opracowała program, który wykorzystuje techniki oceny osobowości i predyspozycji oraz metody uczenia maszynowego do zapewnienia spersonalizowanych rekomendacji kierunków studiów na Politechnice Łódzkiej. W tym celu wybrane zostały algorytmy filtrowania kolaboratywnego oraz \textit{LambdaRank}, powszechnie stosowane w systemach rekomendacji. Rozwiązanie zostało zaimplementowane jako aplikacja internetowa w formie quizu wielokrotnego wyboru, umożliwiając kandydatom otrzymanie dopasowanych rekomendacji na podstawie cech podobnych do aktualnych studentów. Aplikacja zbiera i przetwarza dane, aby wygenerować uporządkowaną listę proponowanych kierunków studiów, ułatwiając dostęp do istotnych informacji bez konieczności bezpośredniego kontaktu ze studentami. Ostatecznie przeanalizowano możliwe usprawnienia i strategie optymalizacji zaproponowanego rozwiązania, w tym ulepszenie zestawu pytań, rozszerzenie zbioru danych wykorzystywanego do trenowania modelu oraz zbadanie alternatywnych algorytmów rekomendacji. 

\bigskip
\section*{Abstract}
\label{chap:StreszczenieEng}
% \indent With the growing number of academic programs offered by universities, prospective students increasingly struggle to make well-informed decisions about their field of study. To address this issue, the author developed a program that utilizes personality and aptitude assessment techniques, along with machine learning methods, to provide personalized study program recommendations at Lodz University of Technology. Collaborative filtering algorithms and \textit{LambdaRank}, commonly used in recommendation systems, were selected for this purpose. The solution is~implemented as a web application in the form of a multiple-choice quiz, allowing candidates to~receive tailored recommendations based on characteristics similar to those of current students. The application collects and processes data to generate ranked study program suggestions, improving accessibility to relevant information without requiring direct interaction with students. Finally, potential improvements and optimization strategies for the proposed solution were analyzed, including refining the question set, expanding the dataset used for model training, and exploring alternative recommendation algorithms.
\bigskip

\textbf{Słowa kluczowe}: \quad automatyzacja  \quad planowanie zadań \quad uczenie przez wzmacnianie \quad algorytm genetyczny \quad produktywność \quad rytm dobowy \quad personalizacja

\noindent\textbf{Keywords}: \quad automation  \quad task planning \quad reinforcement learning \quad genetic algorithm \quad productivity \quad circadian rhythm \quad personalization
